\section{Implementation}

We implemented the design by restructuring \axvisor{}, originally a hypervisor
built on the \arceos{} substrate, into the portable-core-plus-adaptation-layer
form and then adding two new substrates: \asterinas{} and Linux. This section
describes the core restructuring (\S\ref{sec:core-restructuring}), the three
substrate adaptation layers and the semantic mismatches they resolve
(\S\ref{sec:adaptation-layers}), and the KVM-ABI shim
(\S\ref{sec:kvm-shim}). We use Rust throughout, inheriting \axvisor{}'s and
\arceos{}'s language choice, which also lets the contract be expressed as a Rust
trait that each adaptation layer must implement.

\subsection{Core Restructuring}
\label{sec:core-restructuring}

The original \axvisor{} interleaved virtualization logic with direct calls to
\arceos{} services: memory allocation went through \arceos{}'s frame allocator,
vCPU threads were \arceos{} tasks, and interrupt registration used \arceos{}'s
IRQ framework. To make the core portable, we performed a dependency-inversion
refactoring. We defined the substrate interface contract of
Section~\ref{sec:contract} as a trait object that the core holds by reference.
Every former direct call into \arceos{} was replaced by a call through this
trait. Code that was genuinely about virtualization---the vCPU state machine,
stage-2 page-table management, the trap dispatcher, hypercall handlers, and the
virtio device models---remained in the core unchanged in logic. Code that was
about obtaining a resource was moved behind the trait.

The restructuring exposed a useful invariant: the operations the core needs
cluster into the seven service groups of Table~\ref{tab:contract}, and no group
requires more than a handful of operations. After refactoring, the core contains
no substrate-specific symbols; it is linked against whichever adaptation layer
the build selects. This matches goal G1: portability is measured by the thinness
of the adaptation layer, not by conditional compilation in the core.

\subsection{Substrate Adaptation Layers}
\label{sec:adaptation-layers}

We implemented three adaptation layers, each exercising a different point in the
OS design space.

\paragraph{\arceos{} substrate}
\arceos{} is a modular, unikernel-style OS base in which the hypervisor has
historically had direct, low-overhead access to hardware and system services.
Its adaptation layer is the most direct: \arceos{} already exposes physical
frame allocation, task creation, and IRQ registration in forms close to what the
contract requires. Most contract operations map to existing \arceos{} APIs with
light wrapping. This substrate serves as the baseline: it represents the case
where the substrate is already ``hypervisor-friendly'' and imposes minimal
adaptation cost. The main work was re-expressing \arceos{} calls through the
trait so that the core no longer imported \arceos{} symbols directly.

\paragraph{\asterinas{} substrate}
\asterinas{} is a Linux-ABI-compatible, Rust-based framekernel OS with a small,
sound trusted computing base~\cite{peng2025asterinas, peng2024framekernel}. Its
kernel architecture differs substantially from \arceos{}'s unikernel model: it
enforces intra-kernel privilege separation, manages address spaces through a
distinct abstraction, and exposes system services through its own internal
interfaces rather than through the flat, single-address-space model \arceos{}
provides. The \asterinas{} adaptation layer therefore had to reconcile several
semantic mismatches.

The first mismatch concerns physical memory. \arceos{} readily hands out
physical frames; \asterinas{}'s memory subsystem is structured around its own
kernel address-space and safety model, so obtaining raw physical frames and
installing stage-2 mappings required bridging between \asterinas{}'s safe
abstractions and the physical addresses the contract demands. The adaptation
layer recovers physical frame numbers from \asterinas{}'s allocation primitives
and translates contract \code{map\_stage2} calls into \asterinas{}'s page-table
manipulation facilities. The second mismatch concerns execution contexts:
\asterinas{}'s scheduling and thread model differs from \arceos{} tasks, so
\code{vcpu\_run} is implemented on top of \asterinas{} threads that enter the
virtualization mode. The third concerns interrupts: \asterinas{} routes
interrupts through its own framework, and the adaptation layer registers
contract IRQ handlers within that framework while preserving the semantics the
core expects. Despite these mismatches, the adaptation layer is confined to the
contract boundary; the core's virtualization logic required no modification.

\paragraph{Linux substrate}
Linux represents the opposite extreme from \arceos{}: a mature monolithic kernel
with extensive facilities but also strong internal boundaries, a module model, a
privilege structure, and abstractions that often hide the very physical
resources a hypervisor needs. Implementing the contract on Linux is therefore an
instructive stress test: if the contract can be satisfied even where the OS
actively abstracts away physical resources, it is likely sufficient for less
constrained substrates.

The Linux adaptation layer is realized as a kernel module. Several contract
operations required bridging deliberate Linux abstractions. Physical memory
allocation uses contiguous-frame allocators rather than the default page
allocator, because the core needs physical addresses and contiguity for stage-2
mappings. Stage-2 (EPT/NPT) management requires control of the hardware
virtualization structures that Linux normally governs through KVM; the
adaptation layer provisions these without conflicting with the host's own KVM
when co-resident. Interrupt registration bridges to Linux's IRQ subsystem, with
care taken for context semantics: Linux interrupt handlers may run in contexts
where the core's synchronization assumptions need adaptation. The vCPU run loop
is backed by a kernel thread that enters VMX root (x86) or EL2 (Arm). The IOMMU
operations map onto Linux's IOMMU/DMA APIs where available. The key observation
is that Linux \emph{can} provide every contract capability, but each requires
explicitly escaping abstractions that application-facing Linux code never
touches---which is precisely why a hypervisor-grade contract, rather than a
generic OS API, is necessary.

\begin{table}[t]
  \centering
  \caption{Implementation summary across substrates. LoC figures are
  adaptation-layer sizes excluding the shared core. \fillme{Fill with measured
  values}.}
  \label{tab:impl-summary}
  \begin{tabular}{@{}llccl@{}}
    \toprule
    \textbf{Substrate} & \textbf{Model} & \textbf{Adapt. LoC} &
      \textbf{Core reuse} & \textbf{Guest} \\
    \midrule
    \arceos{}    & unikernel   & \fillme{A} & \fillme{100\%} & Linux (KVM ABI) \\
    \asterinas{} & framekernel & \fillme{B} & \fillme{100\%} & Linux (KVM ABI) \\
    Linux        & monolithic  & \fillme{C} & \fillme{100\%} & Linux (KVM ABI) \\
    \bottomrule
  \end{tabular}
\end{table}

\subsection{KVM-ABI Shim}
\label{sec:kvm-shim}

The KVM-ABI shim lives in the portable core and implements the standard ioctl
surface (\code{KVM\_CREATE\_VM}, \code{KVM\_CREATE\_VCPU}, \code{KVM\_RUN},
memory-slot management, irqchip and interrupt routing, and the one-reg register
accessors). Each ioctl translates to core operations that themselves call only
through the substrate contract. Because the shim is substrate-neutral, the same
guest images and the same VMM logic run on all three substrates. In our
evaluation (\S\ref{sec:evaluation}), an unmodified bare-metal Linux kernel boots
as a guest on both the \asterinas{} and Linux substrates, exercising the shim's
create/run/memory/irq paths end to end.

The shim also demonstrates the value of the double-decoupling. Without it, each
substrate would need its own guest-facing interface, and each guest would need
per-substrate adaptation. With it, the guest $\times$ substrate matrix is
covered by one core plus $s$ adaptation layers plus one ABI, rather than
$g \times s$ bespoke paths. Table~\ref{tab:impl-summary} summarizes the
implementation.
